%% ============================================================
%%  chapters/03_interview_prep.tex
%%  Chapter 3: Interview Preparation (MTech, MS, PSU)
%% ============================================================

\chapter{Interview Preparation}
\label{chap:interview-prep}

%% ─────────────────────────────────────────────────────────
\section{General Strategy}
\label{sec:interview-strategy}

\begin{protip}[Core Principle]
  \GATE preparation covers the majority of the technical syllabus. Shift your
  focus from \emph{formula memorization} to \emph{conceptual understanding}---IISc
  panels in particular emphasize fundamentals over rote recall.
\end{protip}

\subsection{For M.Tech Interviews (IITs)}

\begin{pfChecklist}
  \item Master EC core subjects: \textbf{Analog} and \textbf{Digital circuits}.
  \item Do not neglect \textbf{Communications}---IIT panels ask from it
        occasionally.
  \item For AI/ML branches: basic \textbf{C coding}, \textbf{linear algebra},
        and \textbf{matrices} are frequently tested.
  \item Focus on \textbf{fundamentals}, not just formulae.
\end{pfChecklist}

\subsection{For M.Tech Research / MS RA Interviews}

\begin{pfChecklist}
  \item Prepare your UG final year project and any published papers thoroughly.
        Interviews often involve a discussion of your own work.
  \item Expect follow-up questions on the \textbf{technology behind your project}.
  \item IISc panels are particularly rigorous---expect whiteboard problem-solving.
\end{pfChecklist}

\subsection{For PSU Interviews}

See Chapter~\ref{chap:psu}, Section~\ref{sec:psu-prep} for a detailed strategy.

%% ─────────────────────────────────────────────────────────
\section{Branch-wise Focus Areas}
\label{sec:branch-focus}

\begin{tblr}{
    colspec = {X[2,l] X[5,l]},
    hlines, vlines,
    row{1} = {bg=pfBlue, fg=white, font=\sffamily\bfseries},
    row{2,4,6,8} = {bg=pfGrey},
  }
  \toprule
  Branch / Domain & Key Topics to Prepare \\
  \midrule
  VLSI / Microelectronics (ECE)
    & Analog circuits, Digital circuits, Device physics,
      CMOS design, Transistor-level design \\
  Signal Processing / Communications
    & DSP fundamentals, Modulation, Fourier analysis,
      Filter design \\
  Power Electronics / EE (IITs)
    & Power system protection, Stability analysis,
      Electrical machines, Power electronics \\
  AI/ML (Interdisciplinary)
    & Linear algebra, Probability, Basic C/Python,
      Machine learning fundamentals \\
  IISc (all branches)
    & Heavy emphasis on fundamentals; expect whiteboard derivations
      and follow-up questions on incorrect answers \\
  \bottomrule
\end{tblr}

%% ─────────────────────────────────────────────────────────
\section{Recommended Preparation Resources}
\label{sec:resources}

\begin{tblr}{
    colspec = {X[2,l] X[3,l] X[4,l]},
    hlines, vlines,
    row{1} = {bg=pfBlue, fg=white, font=\sffamily\bfseries},
    row{2,4,6} = {bg=pfGrey},
  }
  \toprule
  Subject & Resource & Link / Notes \\
  \midrule
  All Subjects (Interview)
    & PrepFusion Free Course
    & \url{https://prepfusion.in/new-courses/46/content} \\
  Electrical Networks \& RC Circuits
    & NPTEL -- Dr.\ Nagendra Krishnapura
    & Available on YouTube / NPTEL website \\
  Analog Circuits (BJT, MOS, Diode)
    & NPTEL -- Dr.\ Nagendra Krishnapura
    & Same playlist as Networks \\
  Digital Circuits \& Math
    & GATE Syllabus + Standard textbooks
    & Sufficient for most M.Tech interviews \\
  IISc Interview Experiences
    & PrepFusion Interview Prep Video
    & \url{https://youtu.be/gjFeuFKBALQ} \\
  \bottomrule
\end{tblr}

\begin{protip}[Free vs.\ Paid Courses]
  PrepFusion's free courses cover all fundamental concepts required for M.Tech
  interviews. Premium options offer additional technical and aptitude-based
  training. The free course is an excellent starting point before investing in
  paid resources.
\end{protip}

%% ─────────────────────────────────────────────────────────
\section{Preparing for the IAP Branch}
\label{sec:iap}

For the \textbf{Integrated Analog Pathfinder (IAP)} branch or similar
analog-focused programs at IISc:

\begin{pfChecklist}
  \item Study \textbf{oscillator design} thoroughly (Wien Bridge, Colpitts, etc.).
  \item Understand \textbf{negative feedback} deeply---its effect on gain,
        bandwidth, and stability.
  \item Practice \textbf{op-amp circuit} analysis on paper/whiteboard.
  \item Revise \textbf{differential amplifiers}: small-signal model, CMRR,
        differential gain.
  \item Brush up on \textbf{control systems} fundamentals (dominant pole,
        pole-zero plots, step response).
\end{pfChecklist}

%% ─────────────────────────────────────────────────────────
\section{Frequently Asked Questions}
\label{sec:interview-faq}

\begin{faqentry}{What kind of questions are asked in interviews for VLSI and
  related branches?}
  It depends on the panel. Generally, questions come from \textbf{analog} and
  \textbf{digital circuits}. Device physics questions may also be asked to EC
  students. IISc panels additionally ask \textbf{power electronics} questions of
  EE students and occasionally \textbf{communications} questions of EC students.
\end{faqentry}

\begin{faqentry}{How to prepare for MS and M.Tech RA program interviews?}
  Similar to M.Tech interviews, but with an additional component: the
  \textbf{projects and papers} you mentioned in your application are actively
  discussed. Prepare to defend your research interests and explain technical
  choices you made during your UG project.
\end{faqentry}
